\cotefile{157-2}
\scanpage{353}
\mspage{92}

en \deleted{\unsure{quelque|général}}, sous réserve de traiter les 4 cas : \unsure{irréductibles}.

A) Cas $\deleted{\unsure{T}}\tau=\Delta_{\varepsilon}^{2}$, $\tau'=\Delta_{\varepsilon}^{2}$ \inserted{ou $\Delta_{\varepsilon}$}. \deleted{\unsure{Alors la cat.}}

La catégorie $\mathrm{Hom}(\Delta_{\varepsilon}^{2},\Delta_{\varepsilon})$ a deux éléments, \deleted{les} morphismes $\Delta_{\varepsilon}^{2}\longrightarrow\Delta_{\varepsilon}$ correspondant aux \deleted{\unsure{deux}} flèches de transition de $\Delta_{\varepsilon}^{2}$ ($=\mathbb{\Delta}_{2}^{+}$ ou $\mathbb{\Delta}_{2}^{-}$). Dans le cas $\deleted{\unsure{2}}\varepsilon=+1$, ce sont les deux applications
\diagram{\begin{tikzcd}[ampersand replacement=\&,column sep=small,row sep=small]
0\arrow[r]\&1\arrow[r,no head]\&2\\
\unsure{\text{a.}}\&0\arrow[r]\&1
\end{tikzcd}}
de $\{0,1,2\}$ dans $\{0,1\}$, qui appliquent $0$ en $0$, $2$ en $1$ et $1$ soit en $\unsure{1|0}$, \deleted{\unsure{soit en}} \inserted{$(F_0)$} soit en $1$ $(F_1)$. On a évidemment une flèche $F_1\longrightarrow F_0$ (le voyageur pressé \unsure{arrive|arriver} plus vite \unsure{...}). \deleted{\unsure{Donc} $\underline{\mathrm{Hom}}$} \inserted{D'une façon} \unsure{pour} $\varepsilon=-1$
\[
\underline{\mathrm{Hom}}(\Delta_{\varepsilon}^{2},\Delta_{\varepsilon})\simeq\mathbb{\Delta}^{1},\qquad \text{catégorie 0-connexe.}
\]

La catégorie $\mathrm{Hom}(\Delta_{\varepsilon}^{2},\Delta_{\varepsilon}^{2})$ a \unsure{trois} éléments : \deleted{\unsure{resp.}} \unsure{s'identifiant|s'identifient} à l'une des parties non \unsure{vide|vides} de l'ens $\mathrm{Tr}(\tau)$ à deux éléments des flèches de transition $\underline u_1,\underline u_2$. de $\tau$ : il y a $F_{\underline u_1}=F_0$, $\underline F_{u_2}$, $F_{\underline u_1,\underline u_2}$. Si $\varepsilon=+1$, il s'agit des applications croissantes (\unsure{sans|sens} plus) de $\{0,1,2\}$ dans $\{0,1,2\}$, appliquant $0$ \inserted{\unsure{?}} en $0$, $2$ en $2$, et $1$ en ce qu'elle voudra. Donc
\diagram{$\begin{array}{ccl}
1\mapsto2&\text{correspond :}&F_0\\
1\mathrel{\unsure{\mapsto|-}}1&\rule{2cm}{.35pt}&F_1\\
1\mapsto0&\rule{2cm}{.35pt}&F_2
\end{array}$}
et on a bien sûr
\diagram{\begin{tikzcd}[ampersand replacement=\&,column sep=small]
F_2\arrow[r,no head]\& F_1\arrow[r,no head]\& F_0
\end{tikzcd}}
\[
\underline{\mathrm{Hom}}(\Delta_{\varepsilon}^{2},\underline{\Delta}_{\varepsilon}^{2})\simeq\mathbb{\Delta}^{2},
\]
catégorie 0-connexe.

Dans ces deux cas, deux éléments de $\underline{\mathrm{Hom}}(\tau,\tau')$ sont reliés par une homotopie \uline{simple}.

% archive sheet number: 352
\scanpage{354}
\mspage{93}

B) Cas $\tau=\Delta_\varepsilon\Delta_{-\varepsilon}^{\ell}\Delta_\varepsilon$, $\tau'=\mathbb{\Delta}_\varepsilon^2$ ou $\mathbb{\Delta}_\varepsilon$.

Notons que l'on a, pour $\tau,\tau'\unsure{\in|\vdash}\underline{\mathrm{Ty}}$ \unsure{quelconque|quelconques}, un isom. de catégories
\[
\underline{\mathrm{Hom}}(\tau,\tau')\longrightarrow\underline{\mathrm{Hom}}(\tau'{}^\circ,\tau^\circ),
\]
exprimant simplement que
\[
\tau\longmapsto\tau^\circ
\]
est un foncteur \deleted{\unsure{et}} de $\underline{\mathrm{Ty}}$ dans $\underline{\mathrm{Ty}}$, et $\unsure{=|\simeq}$ un automorphisme de cette catégorie. Donc dire que $\underline{\mathrm{Hom}}(\tau,\tau')$ est connexe revient à dire que $\underline{\mathrm{Hom}}(\tau^\circ,\tau'{}^\circ)$ l'est. Donc on peut supposer \unsure{ci-dessous|ci-dessus} que $\varepsilon=+1$ (et pareil dans A !).

\diagram{\begin{tikzcd}[ampersand replacement=\&,column sep=.65em]
\tau:\&{}\arrow[r]\&{}\&{}\arrow[l]\&{}\arrow[l]\&\cdots\&{}\arrow[l]\&{}\arrow[r]\&{}
\arrow[from=1-3,to=1-7,phantom,"{\underbrace{\hspace{3.1cm}}_{\ell}}"']
\end{tikzcd}}
\diagram{\begin{tikzcd}[ampersand replacement=\&,column sep=.8em]
\deleted{\unsure{I}}\tau':\&{}\arrow[r]\&{}\arrow[r]\&{}\&\text{ou}\&{}\arrow[r]\&{}
\end{tikzcd}}

Mais tous les \deleted{foncteurs} morphismes $\tau\mathrel{\unsure{-|\to}}\tau'$ sont triviaux sur l'intervalle critique $I_2$ (rétrograde), ce qui implique que
\[
\underline{\mathrm{Hom}}(\tau,\tau')\xleftarrow{\sim}\underline{\mathrm{Hom}}(\Delta_\varepsilon^2,\tau'),
\]
\deleted{\unsure{en quoi}} \inserted{via} le foncteur de contraction canonique
\[
I_{\unsure{\tau|I}}\mathrel{\unsure{-|\to}}I_{\Delta_\varepsilon^2}=\mathbb{\Delta}^{2}.
\]

Or dans A on a vu que $\underline{\mathrm{Hom}}(\Delta_\varepsilon^2,\tau')$ est connexe (pour $\tau'$ de la forme $\Delta_\varepsilon$ ou $\Delta_\varepsilon^2$).

On a gagné !

% archive sheet number: 353
\scanpage{355}
\mspage{94}

\uline{9.} \fbox{\parbox{.56\linewidth}{\uline{Théorème d'homotopie pour les catégories de chemins}}}

Soient $X$ une catégorie, \inserted{$\Sigma\subset\mathrm{Fl}(X)$ ens. admissible de flèches de $X$} $a,b\in\mathrm{Ob}\,X$. Considérons \deleted{\unsure{les}} pour $\tau\unsure{\in|\vdash}\mathrm{Ty}$ variable, les catégories
\[
\underline{\mathrm{Ch}}_{\Sigma,\tau}(X),\qquad \underline{\mathrm{Ch}}_{\Sigma,\tau}(X;a,b),
\]
qui forment les fibres des catégories Cat.-fibrées
\diagram{\begin{tikzcd}[ampersand replacement=\&,column sep=1em]
\underline{\mathrm{Ch}}_{\Sigma}(X)\arrow[dr,no head,bend right=7]\&\&\underline{\mathrm{Ch}}_{\Sigma}(X;a,b)\arrow[ll,"\mathrm{incl}"']\arrow[dl,no head,bend left=7]\\
\&\underline{\mathrm{Ty}}\&
\end{tikzcd}}
sur la catégorie $\underline{\mathrm{Ty}}$. Ainsi, pour $\tau,\tau'\in\mathrm{Ty}$, et pour \deleted{\unsure{une}}
\[
\varphi:\tau\mathrel{\unsure{-|\to}}\tau',
\]
on a un carré commutatif d'homomorphismes
\diagram{\begin{tikzcd}[ampersand replacement=\&,column sep=large]
\underline{\mathrm{Ch}}_{\Sigma,\tau\unsure{,|'}}(X)\arrow[d,"\varphi_X^*"']\&\underline{\mathrm{Ch}}_{\Sigma,\tau'}(X;a,b)\arrow[l,"\mathrm{incl}"']\arrow[d,"\varphi_{a,b}^*"]\\
\underline{\mathrm{Ch}}_{\Sigma,\tau}(X)\&\underline{\mathrm{Ch}}_{\Sigma,\tau}(X;a,b)\arrow[l,"\mathrm{incl.}"']
\end{tikzcd}}\hfill,
avec transitivité pour $\varphi$ variable (composition $\psi\circ\varphi$ de flèches dans $\underline{\mathrm{Ty}}$).

En fait, pour $\varphi$ variable dans $\underline{\mathrm{Hom}}(\tau,\tau')$, les foncteurs images inverses $\varphi_X^*,\varphi_{\unsure{ab|a,b}}^*$ dépendent \uline{fonctoriellement} de $\varphi$ ; i.e. on a des \uline{foncteurs}
\diagram{\begin{tikzcd}[ampersand replacement=\&,column sep=small,row sep=small]
(*)\&\underline{\mathrm{Hom}}(\tau,\tau')\arrow[r]\arrow[dr,bend right=15]\&\underline{\mathrm{Hom}}\bigl(\underline{\mathrm{Ch}}_{\Sigma\unsure{'|{}},\tau'}(X),\underline{\mathrm{Ch}}_{\Sigma,\tau\deleted{\unsure{'|*}}}(X)\bigr)\\
\&\&\underline{\mathrm{Hom}}\bigl(\underline{\mathrm{Ch}}_{\Sigma,\tau'}(X\unsure{,|;}a,\deleted{\unsure{=|-}}\unsure{b|t}),\underline{\mathrm{Ch}}_{\Sigma,\tau}(X;a,b)\bigr).
\end{tikzcd}}

% archive sheet number: 354
\scanpage{356}
\mspage{95}

Cela provient de l'inclusion
\[
\underline{\mathrm{Hom}}(\tau,\tau')\subset\underline{\mathrm{Hom}}^{\varepsilon}(I_\tau,I_{\tau'})
\]
et de la description des catégories $\underline{\mathrm{Ch}}_\tau,\underline{\mathrm{Ch}}_{\tau'}$ comme catégories de foncteurs, telle
\[
\underline{\mathrm{Ch}}_{\Sigma,\tau}(X)\simeq\underline{\mathrm{Hom}}_{\Sigma_\tau,\Sigma}(I_\tau,X),
\]
où $\Sigma_\tau$ désigne l'ens. des flèches dans $I_\tau$ qui sont soit \deleted{\unsure{des iso}} inversibles, soit \uline{rétrogrades}. (Cet ens. de flèches est admissible, car stable par composition.) On a désigné par $\underline{\mathrm{Hom}}_{\Sigma_\tau,\Sigma}$ la sous-catégorie pleine de $\underline{\mathrm{Hom}}$ formée des foncteurs qui appliquent $\Sigma_\tau$ dans $\Sigma$. Or de façon générale, si
\[
(I,\Sigma_I)(I',\Sigma_{I'}),\quad(X,\Sigma)
\]
sont trois catégories, munies d'ens. admissibles de flèches, \deleted{\unsure{on trouve les foncteurs suivants}} \inserted{considérant les foncteurs} $\varphi:I\mathrel{\unsure{-|\to}}I'$ \deleted{\unsure{un foncteur}} avec $\varphi(\Sigma_I)\subset\Sigma_{I'}$, alors le foncteur \deleted{\unsure{f}} \unsure{usuel|usuelle|naturel}
\[
\underline{\mathrm{Hom}}(I,I')\longrightarrow\underline{\mathrm{Hom}}\bigl(\underline{\mathrm{Hom}}(I',X),\underline{\mathrm{Hom}}(I,X)\bigr)
\]
(transposé du foncteur de composition
\[
\begin{aligned}
\underline{\mathrm{Hom}}(I,I')\times\underline{\mathrm{Hom}}(I',X)&\longrightarrow\mathrm{Hom}(I,X))\\
(\varphi,F')&\longmapsto F'\circ\varphi),
\end{aligned}
\]
\deleted{\unsure{induit}} \inserted{se généralise en} un foncteur
\[
\underline{\mathrm{Hom}}_{\Sigma_I,\Sigma_{I'}}(I,I')\longrightarrow
\underline{\mathrm{Hom}}\bigl(\underline{\mathrm{Hom}}_{\Sigma_{I'},\Sigma}(I',X),\underline{\mathrm{Hom}}_{\Sigma_I,\Sigma}(I,X)\bigr),
\tag*{($\unsure{*|**}$)}
\]
transposé du foncteur de composition
\[
\underline{\mathrm{Hom}}_{\Sigma_I,\Sigma_{I'}}(I,I')\times\underline{\mathrm{Hom}}_{\Sigma_I,\Sigma}(I',X)\mathrel{\unsure{-|\to}}\underline{\mathrm{Hom}}_{\Sigma_I,\Sigma}(I,X)
\]
induit par $(*)$. \deleted{\unsure{Comme tout}}

% archive sheet number: 355
\scanpage{357}
\mspage{96}

Comme \deleted{\unsure{toute}}
\[
\varphi:\tau\mathrel{\unsure{-|\to}}\tau'
\]
induit
\[
\varphi:I_\tau\mathrel{-}I_{\tau'}
\]
transformant flèches rétrogrades en flèches identiques ou \unsure{régtrogrades|rétrogrades}, i.e.
\[
\varphi(\Sigma_\tau)\subset\Sigma_{\tau'},
\]
on trouve le premier des deux foncteurs $(*)$ annoncés p. 94, qui est associé par transposition au foncteur $\gamma$
\diagram{\begin{tikzcd}[ampersand replacement=\&,column sep=small,row sep=large]
(*)\&\underline{\mathrm{Hom}}(\tau,\tau')\times
\underbrace{\underline{\mathrm{Ch}}_{\Sigma,\tau'}(X)}_{\simeq\underline{\mathrm{Hom}}_{\Sigma_{\tau'},\Sigma}(I_{\tau'},X)}
\arrow[d,"\mathrm{can}"']\arrow[r,"\gamma"]\&
\underline{\mathrm{Ch}}_{\Sigma,\tau}(X)\arrow[d,phantom,"\parallel"]\\
\&\underline{\mathrm{Hom}}_{\Sigma_\tau,\Sigma_{\tau'}}(I_\tau,I_{\tau'})\times\underline{\mathrm{Hom}}_{\Sigma_{\tau'},\Sigma}(I_{\tau'},X)
\arrow[r,dotted,"\gamma_0"]\&\underline{\mathrm{Hom}}_{\Sigma_\tau,\Sigma}(I_\tau,X)
\end{tikzcd}}
\unsure{déduit} du foncteur de composition $\gamma_0$ dans ce diagramme : $\gamma=\gamma_0\circ\mathrm{can}$. D'autre part, ce foncteur de composition $\gamma$ induit visiblement un foncteur sur les sous-catégories de chemins à \deleted{\unsure{origine}} \inserted{source} et \deleted{\unsure{extrémité}} \inserted{but} fixés
\[
\underline{\mathrm{Hom}}(\tau,\tau')\times\underline{\mathrm{Ch}}_{\Sigma,\tau'}(X;a,b)\mathrel{-}\underline{\mathrm{Ch}}_{\Sigma,\tau}(X;a,b),
\]
d'où la deuxième des deux flèches annoncées dans $(*)$ p. 94.

Cela montre donc que si
\diagram{\begin{tikzcd}[ampersand replacement=\&,column sep=small]
\varphi_1,\varphi_2:\&\tau\arrow[r,shift left=.5ex,no head]\arrow[r,shift right=.5ex,no head]\&\tau'
\end{tikzcd}}
sont deux morphismes de $\tau$ dans $\tau'$ homotopes dans $\underline{\mathrm{Hom}}(\tau,\tau')$, alors les \deleted{deux} foncteurs induits entre \deleted{\unsure{chemins}} \inserted{cat. de} $\Sigma$-chemins ou \inserted{entre cat. de} $\Sigma$-chemins à \unsure{extrémités} \unsure{précisées|précisés} :

% archive sheet number: 356
\scanpage{358}
\mspage{97}

\diagram{\begin{tikzcd}[ampersand replacement=\&,column sep=small]
\varphi_{1X}^*,\varphi_{2X}^*:\&\deleted{\unsure{\mathrm{Hom}}}\underline{\mathrm{Ch}}_{\Sigma,\tau'}(X)\arrow[r,shift left=.5ex,no head]\arrow[r,shift right=.5ex,no head]\&\underline{\mathrm{Ch}}_{\Sigma,\tau}(X),
\end{tikzcd}}
\deleted{\unsure{et pour les}}
\[
(\varphi_1)_{a,b}^*,(\varphi_2)_{a,b}^*:\underline{\mathrm{Ch}}_{\Sigma,\tau'}(X;a,b)\mathrel{\unsure{-|\to}}\underline{\mathrm{Ch}}_{\Sigma,\tau}(X;a,b)
\]
sont \inserted{également} homotopes. \deleted{\unsure{Or}}

Or le lemme d'homotopie nous assure que l'hyp. $\varphi_1\sim\varphi_2$ est \uline{toujours} satisfaite. Donc

\uline{Théorème} \inserted{\uline{d'homotopie \unsure{sur} les opérations \inserted{universelles} sur les chemins.}} Soient $\deleted{\unsure{\{} }X,\Sigma,a,b$ comme dessus, et $\tau,\tau'\in\mathrm{Ty}$ telles que $\mathrm{Hom}(\tau,\tau')\ne\varnothing$. Alors les foncteurs $\varphi_X^*$, $\varphi_{X,a,b}^*$ induits par \unsure{une|un} $\varphi\in\mathrm{Hom}(\tau,\tau')$ :
\[
\left\{\begin{aligned}
\varphi_X^* &: \underline{\mathrm{Ch}}_{\Sigma,\tau'}(X)\longrightarrow\underline{\mathrm{Ch}}_{\Sigma,\tau}(X)\\
\varphi_{X,a,b}^* &: \underline{\mathrm{Ch}}_{\Sigma,\tau'}(X;a,b)\mathrel{-}\underline{\mathrm{Ch}}_{\Sigma,\tau}(X;a,b)
\end{aligned}\right.
\]
ne dépendent pas de $\varphi$, \uline{à homotopie près}.

On peut préciser cet \deleted{\unsure{é}} \unsure{énoncé|l'énoncé}, de telle façon que l'énoncé pour les $\varphi_{X;ab}^*$ apparaisse, comme il devrait, comme corollaire pour celui des $\varphi_X^*$. Pour ceci, on va regarder les catégories de chemins \unsure{sur|dans} $X$ comme \inserted{cat.} \uline{au dessus de $X$}, par les foncteurs ``source et but'' :
\diagram{\begin{tikzcd}[ampersand replacement=\&]
\underline{\mathrm{Ch}}_{\Sigma,\tau}(X)\arrow[dr,"{(\sigma_\tau,\beta_\tau)}"']\&\&\underline{\mathrm{Ch}}_{\Sigma,\tau'}(X)\arrow[dl,"{(\sigma_{\tau'},\beta_{\tau'})}"]\\
\&X\times X\&
\end{tikzcd}}

% archive sheet number: 357
\scanpage{359}
\mspage{98}

Alors le foncteur $(\unsure{\alpha|*})$ (p. 94) \deleted{\unsure{peut}} se factorise par la sous-catégorie $\underline{\mathrm{Hom}}_{\mathrm{Cat}/X\times X}(\,.\,,\,.\,)$ de $\underline{\mathrm{Hom}}_{\mathrm{Cat}}(\,.\,,\,.\,)$, formée des foncteurs qui commutent aux foncteurs de projection dans $X\times X$, et les hom. entre tels foncteurs qui s'effectuent par des flèches ``verticales'' \unsure{de|dans} la catégorie but (i.e. par des flèches au dessus des flèches identiques de la catégorie base $X\times X$).
\[
\underline{\mathrm{Hom}}(\tau,\tau')\longrightarrow\underline{\mathrm{Hom}}_{\mathrm{Cat}/X\times X}
\bigl(\underline{\mathrm{Ch}}_{\Sigma,\tau'}(X),\underline{\mathrm{Ch}}_{\Sigma,\tau}(X)\bigr)
\]
\unsure{:}

Donc dans l'énoncé du théorème d'homotopie, on peut préciser que l'homotopie se fait dans cette sous-catégorie, elle s'effectue par \deleted{\unsure{une|des}}\unsure{des|une} \unsure{suites|suite|bouts} d'homotopies simples qui sont des morphismes de foncteurs entre \deleted{$X\times X$-}catégories au dessus de $X\times X$. Cela implique alors l'énoncé d'homotopie sur les $\varphi_{X,a,b}^*$.

% archive sheet number: 358
\scanpage{360}
\mspage{99}

10. \uline{Questions de comparaison \unsure{pour} les catégories Cat.-\unsure{fibrée|fibrées} \deleted{\unsure{Cat}} satisfaisant au ``th. d'homotopie''} (base \unsure{filtrante ?|filtrante}).

Nous allons généraliser maintenant au cas d'une catégorie (Cat)-fibrée sur une autre
\diagram{\begin{tikzcd}[ampersand replacement=\&]
C\arrow[d,"\unsure{\pi|p}"']\\ T,
\end{tikzcd}}
\margin{\unsure{(}}
en pensant surtout au cas $T=\underline{\mathrm{Ty}}$, $C=\underline{\mathrm{Ch}}_\Sigma(X;a,b)$. L'hypothèse essentielle est la suivante

\margin{(Hypothèse d'homotopie)}Si $t,t'\in\mathrm{Ob}\,T$ sont tels que $\mathrm{Hom}(t,t')\ne\varnothing$, alors les foncteurs $\varphi^*:C_{t'}\mathrel{\unsure{-|\to}}C_t$ associés aux $\varphi\in\mathrm{Hom}(t,t')$ sont homotopes.

\deleted{\unsure{Soit donc}} Ils définissent donc une seule et \unsure{$=$|même} flèche dans $\mathrm{Cat\,hot}$. \deleted{\unsure{So}} Je désigne par $T_{\mathrm{ord}}$ \inserted{ou $T_0$} l'ens $\mathrm{Ob}(T)$, avec comme relation d'ordre celle déduite de la structure de catégorie :
\[
t\leq t'\quad\Longleftrightarrow\quad\mathrm{Hom}_T(t,t')\ne\varnothing,
\]
\deleted{\unsure{donc (ainsi} $T_{\mathrm{ord}}$} de sorte qu'on a un foncteur
\[
T\longrightarrow T_0
\]
qui est \uline{universel} pour les foncteurs de $T$ dans des catégories ``ordonnées'' (i.e. où \unsure{\deleted{deux}|toute} flèche \unsure{est} mono-morphique) \inserted{(ou pour les foncteurs $F$ sur $T$ tels que $F$ soit constant sur tout $\mathrm{Hom}(t,t')$)}. L'\unsure{hypth.|hypothèse|hyp.} d'homotopie peut s'exprimer en disant que le \deleted{\unsure{a lieu de la manière toute}} \inserted{foncteur composé}
\diagram{\begin{tikzcd}[ampersand replacement=\&,column sep=small,row sep=small]
\deleted{\unsure{C}}\,T^\circ\arrow[r]\&\mathrm{Cat}\arrow[r]\&\mathrm{Cat\,hot}\\
t\arrow[r,mapsto]\& C_t\&
\end{tikzcd}}

% archive sheet number: 359
\scanpage{361}
\mspage{100}

se factorise (évidemment de façon unique) en un foncteur
\[
H_C:T_0^\circ\longrightarrow\mathrm{Cat\,hot}
\]
i.e. qu'on a \deleted{\unsure{donc}} un diagramme commutatif
\diagram{\begin{tikzcd}[ampersand replacement=\&]
T^\circ\arrow[r,"{t\mapsto C_t}"]\arrow[d]\&\mathrm{Cat}\arrow[d]\\
T_0^\circ\arrow[r,"{\substack{H_C\\\deleted{\unsure{C}}}}"]\&\mathrm{Cat\,hot}.
\end{tikzcd}}

Nous allons supposer de plus

\margin{(Hypothèse de filtration)}L'ens ordonné $T_0$ associé à $T$ est \uline{filtrant} décroissant \inserted{(filtrant à gauche)}, i.e. si $t,t'\in\mathrm{Ob}\,T$, il existe $t''\in\mathrm{Ob}\,T$ \deleted{\unsure{tel que}} des flèches $t^{\unsure{\prime\prime|\prime\prime\prime}}\mathrel{\unsure{-|\to}}t$, $t''\mathrel{-}t'$.

Cette hypothèse est bien satisfaite dans $\underline{\mathrm{Ty}}$, car pour $\tau,\tau'\in\mathrm{Ob}\,\underline{\mathrm{Ty}}=\mathrm{Ty}$, il suffit de prendre $\tau\tau'$, qui s'envoie dans $\tau$ par $\mathrm{id}_\tau*p_{\tau'}$, dans $\tau'$ par $p_\tau*\mathrm{id}_{\tau'}$, où $p_\tau:\tau\longrightarrow1$ est le morphisme de $\tau$ dans l'objet final. \deleted{\unsure{L|E}}

Ainsi, $\deleted{\unsure{\Phi|H}}H_C$ définit un Ind-objet de $\mathrm{Cat\,hot}$, \inserted{que nous allons désigner par $\underline{\underline{\unsure{\mathcal C|C}}}$}. Si on suppose de plus

\margin{Hypothèse de dénombrabilité (tjrs satisfaite si $\mathrm{card}\,T_0\leq\aleph_0$)}$T_0$ admet un sous-ensemble \uline{dénombrable} cofinal à gauche,

alors cet ind-objet de $\mathrm{Cat\,hot}$ peut être défini, \unsure{\deleted{de}|de} \unsure{beaucoup|\deleted{beaucoup}|\deleted{la manière}} de \unsure{façons|façon}\unsure{,} en termes

% archive sheet number: 360
\scanpage{362}
\mspage{101}

d'un système inductif dans $\mathrm{Cat}$ (pas seulement dans $\mathrm{Cat\,hot}$)
\[
(C_n)_{n\in\mathbb N},
\]
d'où \unsure{on|en} récupère le ind-objet $\underline{\underline C}$ dans $\mathrm{Cat\,hot}$ en passant au syst. inductif correspondant dans $\mathrm{Cat\,hot}$, puis à l'ind-objet associé.

Pour construire le système inductif $C_n$, il suffit de prendre une chaine infinie
\diagram{\begin{tikzcd}[ampersand replacement=\&,column sep=.7em]
t_0\&t_1\arrow[l]\&t_2\arrow[l]\&\cdots\&t_n\arrow[l]\&t_{n+1}\arrow[l]\&\cdots\arrow[l]
\end{tikzcd}}
de morphismes dans $T$, telle que l'ens des $t_i$ soit cofinal \inserted{à gauche} dans $T_0$. \deleted{Ainsi} On va supposer de plus $C$ \uline{strictement fibrée} sur \deleted{$\unsure{\Sigma|I}$} $T$ (i.e., provenant d'un véritable foncteur $T^\circ\longrightarrow\mathrm{Cat}$,) i.e. on peut choisir les foncteurs images inverses \deleted{$\unsure{\varphi}$}
\[
\varphi^*:C_{t'}\mathrel{\unsure{-|\to}}C_t
\]
de façon transitive stricte, pas seulement transitive à iso. près) -- cas auquel on peut tjrs se ramener d'ailleurs par l'astuce de Giraud. Cette hypothèse est d'ailleurs vérifiée dans le cas $\underline{\mathrm{Ch}}_\Sigma(X;a,b)$ sur $\underline{\mathrm{Ty}}$. Ainsi, \deleted{\unsure{les}}
\[
(C_{t_i})_{i\in\mathbb N}
\]
forment un système \deleted{projectif} \inserted{inductif} \unsure{de|des} \unsure{catégorie|catégories}, \unsure{indexés|indexées|indexé} par $\mathbb N$, le morphisme de transition
\[
C_{\unsure{n}}\longrightarrow C_{n'}\qquad\text{i.e.}\quad C_{t_n}\mathrel{-}C_{t_{n'}}
\]

% archive sheet number: 361
\scanpage{363}
\mspage{102}

étant le foncteur image inverse $\varphi_{n,n'}^*$, où
\[
\varphi_{n,n'}:t_{n'}\longrightarrow t_n
\]
est le composé des $t_{n'}\mathrel{\unsure{-|\to}}t_{n'-1}\mathrel{-}\ldots\mathrel{-}t_n$. (Si on n'avait pas transitivité stricte, on définirait $\varphi_{n,n'}^*$ comme le composé des foncteurs image inverse par les flèches de transition successives $t_{n'}\mathrel{-}t_{n'-1}\mathrel{-}\cdots\mathrel{-}t_{n\unsure{-|+}1}\mathrel{-}t_n$, ce ne serait $\varphi_{n,n'}^*$ qu'à isom. can. près.)

Soit \deleted{$C_\infty$}
\[
C_\infty=\underset{n\unsure{.}}{\varinjlim}\,C_n=\underset{n}{\underline{\lim}}\,C_{t_n}
\]
la catégorie $\varinjlim$. (Elle sera surtout sympa si on a pu choisir le système $(t_n)_{n\in\mathbb N}$ de telle façon que les foncteurs de transition
\[
C_{t_n}\mathrel{-}C_{t_{n\unsure{-|+}1}}
\]
soient pleinement fidèles, et peut-être même \unsure{mono|monos}, de telle façon que les $C_n$ s'identifient à des sous-catégories pleines de leur limite inductive $C_\infty$.) Je voudrais comprendre les relations précises entre les trois genres d'objets suivants :

\textcircled{a} $C$ elle-même, \deleted{\unsure{comme}} \inserted{(comme objet de $\mathrm{Cat}$, pas de $\mathrm{Cat}/T$\unsure{.|!})} et surtout les objets de $\mathrm{Cat\,hot}$ et de $\mathrm{Hot}$ qui lui correspondent.

% archive sheet number: 362
\scanpage{364}
\mspage{103}

\textcircled{b} $C_\infty$, et les objets de $\mathrm{Cat\,hot}$ et de $\mathrm{Hot}$ qui lui correspondent. Dans quelle mesure ces derniers sont ils indépendants du choix du syst. proj. cofinal à g. $(t_n)_{n\in\mathbb N}$ dans $T$ ? Quel est le rôle de $C_\infty$ ? Est-elle \inserted{canoniquement} isomorphe à $C$ dans $\mathrm{Cat\,hot}$, ou dans $\mathrm{Hot}$ ? (Auquel cas la question de son ``rôle'' me semblerait résolue !)

\textcircled{c} Le ind-objet $H_C$ dans $\mathrm{Cat\,hot}$, et \unsure{\deleted{le}|le} ind-objet de $\mathrm{Hot}$ qu'il définit.

Je note que le Ind-objet $H_C$, et aussi les catégories $C_\infty$ construites comme dessus, ne changent pas quand on remplace $T$ par une sous-catégorie \inserted{$T'$} cofinale à gauche, ou \deleted{par} aussi par une catégorie de la forme $T_{/t_0}$ \inserted{$(t_0\in\mathrm{Ob}\,T)$}, qui donne lieu à \deleted{\unsure{le} \illegible{1} \unsure{de|à} \unsure{cat.}} \inserted{l'ens. ordonné} \deleted{\unsure{ord}} associé
\[
\bigl(T_{/t}\bigr)_0=T_{0/t}=\{t^{\unsure{2|\tau|\prime|?}}\in T_0\mid \unsure{t_0|t}\leq t_0\}.
\]

Si la relation entre les \deleted{catégories} \inserted{objets} dans a), et ceux dans b) \unsure{ou} dans c), est très étroite, on en conclurait que le ``type'' \deleted{d'homotopie} de $C$ dans $\mathrm{Cat\,hot}$, ou dans $\mathrm{Hot}$ suivant les cas, ne change pas en remplaçant $C$ par

% archive sheet number: 363
\scanpage{365}
\mspage{104}

$C\times_T T'$, ou par $C_{/t_0}\simeq C\times_T T_{/t_0}$. \deleted{\unsure{Dans ce cas}} \inserted{Bien} mieux, elle ne changerait pas en remplaçant $T$ par $(\mathbb N)^\circ$, la catégorie (ordonnée) opposée à $\mathbb N$ (ayant donc un unique objet final), et la catégorie fibrée $C$ sur $T$ par la catégorie $C_{\underline t}$ induite par $C$ sur $(\mathbb N)^\circ$, via le foncteur
\[
\underline t:(\mathbb N)^\circ\mathrel{\unsure{-|\to}}T
\]
défini par le système projectif $\underline t=(t_n)_{n\in\mathbb N}$ dans $T$. Cela signifierait donc que pour comprendre le type homotopique (au sens strict, i.e. dans $\mathrm{Cat\,hot}$, ou large i.e. dans $\mathrm{Hot}$) de $C$, on peut \unsure{le|la} ``nettoyer'' de façon draconienne, la base $T$ étant \inserted{vue comme} une sorte de gangue un peu \unsure{vaseuse|vasarde}, ou des \unsure{scories} \unsure{pléthoriques},

\textcircled{1} De toutes façons, si $T'$ est une \inserted{autre} catégorie filtrante à gauche, et
\[
\unsure{g|\varphi}:T'\longrightarrow T
\]
un foncteur cofinal à gauche, \deleted{\unsure{(de}} de sorte que $C\deleted{\unsure{/T}}$ sur $T$, et $C'=C\times_T T'$ sur $T'$, définissent \inserted{essent.} le ``$=$'' objet de $\mathrm{Ind\,Cat\,hot}$, et que \unsure{la catégorie|les catégories} $C'_\infty$ \inserted{pour $C'$ sur $T'$} \unsure{fait|font} partie des catégories $C_\infty$ pour $C$ sur $T$ -- on peut se demander

% archive sheet number: 364
\scanpage{366}
\mspage{105}

si le foncteur canonique
\[
C'\mathrel{\unsure{-|\to}}C
\]
est \inserted{également tjrs} un iso. dans $\mathrm{Cat\,hot}$, ou sinon, \unsure{du|des} moins dans $\mathrm{Hot}$.

Un intérêt particulier s'attache à $\pi_0(\unsure{\mathcal C|C})$, et de plus à sa variance pour $C\deleted{\unsure{/T}}$ \inserted{sur $T$} variable de la façon suivante. Si $\deleted{\unsure{\Sigma|S}}S$ est dans $\mathrm{Cat}$, on regarde les catégories
\[
\underline{\mathrm{Hom}}_{\mathrm{Cat}}(S,C_t)\qquad(t\in T),
\]
on trouve un foncteur
\[
t\longmapsto\underline{\mathrm{Hom}}_{\mathrm{Cat}}(S,C_t),\qquad T^\circ\mathrel{-}(\mathrm{Cat}),
\]
d'où une catégorie fibrée sur $T$, que je vais désigner par
\[
\underline{\mathscr H\mathrm{om}}_T(S_T,C).
\]

(Sauf erreur, elle s'interprète comme la catégorie fibrée sur $T$ correspondant au champ \inserted{sur $\mathrm{Top}(T)=T^\wedge$} des morphismes de $T^{\unsure{\wedge|\circ}}$/champs de $S_T$ dans $C$, où $S_T=S\times T$. (considéré comme catégorie \inserted{fibrée} sur $T$ à fibre \inserted{\unsure{cst}} \deleted{constant} $\unsure{S|D}$) dans $C$. Ici j'identifie \unsure{catégorie fibrée|catégories fibrées} sur $T$, et champs sur le topos $T^\wedge$.)

Il est immédiat que \deleted{cette} la propriété d'homotopie (p. 99) est stable par passage de $\mathcal C$ à $\underline{\mathscr H\mathrm{om}}_T(S_T,C)$. On s'intéresse à la variation de
\[
\pi_0\bigl(\underline{\mathscr H\mathrm{om}}_T(S_T,\mathcal C)\bigr)
\]

% archive sheet number: 365
\scanpage{367}
\mspage{106}

comme contrafoncteur en $S$,
\[
H:S\longmapsto\pi_0\,\underline{\mathscr H\mathrm{om}}_T(S_T,\unsure{\mathcal C|C})\qquad(\mathrm{Cat})^\circ\longrightarrow(\mathrm{Ens}),
\]
et à une éventuelle expression de ce \unsure{cfoncteur|foncteur} $H$ en termes du ind-objet $H_C$ de $\mathrm{Cat\,hot}$, par la formule hypothétique
\[
H(S)\simeq\mathrm{Hom}_{\inserted{\mathrm{Ind}(\mathrm{Cat\,hot})}}(\gamma(S),H_C),
\]
où
\[
\gamma:\mathrm{Cat}\longrightarrow\mathrm{Cat\,hot}
\]
est le foncteur canonique. Si on avait un tel isomorphisme \inserted{canonique} fonctoriel en $S$, cela donnerait une description (à isom unique près) du ind-objet $H_C$ en termes du foncteur précédent $H$ sur $\mathrm{Cat}$. \deleted{\unsure{I}}

Il me vient un doute ici, en vue de l'exemple $\deleted{\unsure{C}}T=\underline{\mathrm{Ty}}$, $C=\underline{\mathrm{Ch}}_\Sigma(X;a,b)$. Les trois types d'objets a) b) c) ont le $\unsure{=|\text{même}}$ $\pi_0$, et il est sûr que ``c'est le bon''. D'autre part, pour une catégorie \unsure{paramètre|paramétrante} variable $S$, il est sûr aussi que le \unsure{bon} foncteur en $S$ est
\[
S\longmapsto\pi_0\bigl(\underline{\mathrm{Hom}}_T(S_T,\mathcal C)\bigr)\simeq\mathrm{Hom}_{\mathrm{Ind\,Cat\,hot}}(S,H_C)
\]
isomorphe aussi, \uline{si $\deleted{\unsure{D|S}}\unsure{S|D}$ est finie}, à
\[
\deleted{\unsure{\pi_0}}\mathrm{Hom}_{\mathrm{Cat\,hot}}(S,C_\infty),
\]
ce qui montre que, du moins pour des arguments $S\in(\mathrm{Cat})$ qui sont des cat. \uline{finies}, les objets

% archive sheet number: 366
